\subsubsection{问题四的数据接口与有限动作集合}

问题四从问题一的 150 条原始计划重新开始，不读取问题二的工作簿作为约束，也不把问题三新增的 C 类计划混入本问。基准资源域仍为 \([0,643)\times[0,100)\)，事件采用半开区间，所有时间端点和调整量按整数资源单位处理。

每份计划只能选择一种动作。A、B、C 三类均可保持、频率平移、首次时间平移或撤销；C 类额外允许改变重复使用之间的空闲间隔。频率平移量和首次时间平移量分别满足
\[
\delta f\in\{-10,\ldots,-1,1,\ldots,10\},\qquad
\delta t\in\{-5,\ldots,-1,1,\ldots,5\}.
\]
对 C 类计划，令原空闲间隔为 \(g_i=8\)，间隔调整量为
\[
\delta g\in\{-8,\ldots,-1,1,\ldots,10\},
\qquad g_i'=8+\delta g\in\{0,1,\ldots,18\}.
\]
选择间隔动作时，频段区间和首次时间区间保持原值；选择频率或时间平移动作时，\(\delta g=0\)。使用次数、频段宽度和单次占用时长均不改变。

\subsubsection{间隔调整下的重复事件递推}

若计划 \(i\) 的首次时间区间为 \([s_i,e_i)\)，单次占用时长为 \(l_i=e_i-s_i\)，调整后的空闲间隔为 \(g_i'\)，则第 \(k\) 次事件为
\begin{equation}
T_{ik}(g_i')=
\left[s_i+k(l_i+g_i'),\ e_i+k(l_i+g_i')\right),
\qquad k=0,1,\ldots,n_i-1.
\label{eq:q4-event-expansion}
\end{equation}
间隔变化会联动第二次及以后的全部事件，因此每个候选动作先完整展开，再进行边界和资源格检查；不以单个事件平移或平均占用替代完整递推。

\subsubsection{问题四资源格模型与目标}

令 \(\mathcal A_i^{(4)}\) 为计划 \(i\) 的 Q4 合法动作集合，令 \(x_{ia}\in\{0,1\}\) 表示是否选择动作 \(a\)，令 \(\mathcal K_{ia}^{(4)}\) 表示该动作的完整重复事件占用的离散时频资源格集合。每份计划恰选择一种动作：
\begin{equation}
\sum_{a\in\mathcal A_i^{(4)}}x_{ia}=1,
\qquad i\in\mathcal I.
\label{eq:q4-one-action}
\end{equation}
对每个资源格 \(c\)，加入
\begin{equation}
\sum_{(i,a):\,c\in\mathcal K_{ia}^{(4)}}x_{ia}\le 1,
\qquad c\in\mathcal C.
\label{eq:q4-cell-clique}
\end{equation}
该约束保证所有未撤销计划的完整事件在两个维度上均无冲突。基准场景共生成 6\,103 个合法候选动作，其中 5\,953 个为活动候选；按同一计划的外部冲突集合做支配消元后保留 5\,533 个活动候选。420 个被消去动作均有同计划、外部冲突集合为其子集的保留动作，支配替换核验通过。

为保持与问题二相同的优先级口径，Q4 仍按
\[
R\rightarrow R_A\rightarrow R_B\rightarrow M\rightarrow M_A\rightarrow M_B\rightarrow S_{10}^{(4)}
\]
比较，其中前六层分别表示撤销和调整计划数，末级候选指标为
\begin{equation}
S_{10}^{(4)}=\sum_i\left(|\delta f_i|+2|\delta t_i|+|\delta g_i|\right).
\label{eq:q4-shift-score}
\end{equation}
本阶段只闭合决定方案规模的最小撤销层；\(R=3\) 内的部分次级目标仍以条件候选记录，不写成完整字典序最优。

\subsubsection{最小撤销数的上下界证明}

首先由 \texttt{proof\_R\_le\_3.json} 构造 3 台撤销的无冲突方案，并通过连续半开区间、离散资源格和独立重建检查。因此得到可行上界 \(R\le3\)。

为证明 \(R\le2\) 不可行，若存在少于 2 台撤销的可行方案，可以再撤销任意计划补足到恰好 2 台，且撤销不会制造新的资源冲突，故只需考察恰好 2 台撤销。两台撤销在 A、B、C 三类之间只有
\[
(0,0,2),(0,1,1),(0,2,0),(1,0,1),(1,1,0),(2,0,0)
\]
六种类别分拆。对六个分支分别保留全部计划身份和合法活动动作，建立 \(R=2\) 的 SAT 判定模型；六个分支均返回 \texttt{UNSAT}。分支覆盖和支配消元均由独立核验文件检查通过，因此 \(R\le2\) 在本文有限动作集合内不可行。该阈值核验采用 CaDiCaL\cite{cadical}，但其输入候选全集和资源口径与主模型相同。结合 3 台可行见证，得到
\begin{equation}
R_{Q4}^{\ast}=3.
\label{eq:q4-min-revoke}
\end{equation}
这里的结论只针对整数时频资源格、\(H=643\) 和上述有限动作集合，不外推到连续调整或未列出的动作。

\begin{table}[H]
\centering
\caption{问题四最小撤销层的证据状态}
\label{tab:q4-proof}
\begin{tabular}{lccc}
\toprule
检验层 & 方案/阈值 & 证据状态 & 允许的结论 \\
\midrule
可行上界 & \(R\le3\) & 可行且独立验证通过 & 3 台撤销可行 \\
全局下界 & \(R\le2\) & 六个类别分支均 \texttt{UNSAT} & 2 台及以下不可行 \\
最小撤销层 & \(R=3\) & 上下界闭合 & 当前有限模型下最小撤销数为 3 \\
\bottomrule
\end{tabular}
\end{table}

\subsubsection{当前条件候选与工作簿核验}

在固定 \(R=3\)、\(R_A=0\)、\(R_B=1\) 的条件下，当前候选链给出 142 个调整计划，但 \(R_B=1\)、\(M=142\) 和 \(S_{10}^{(4)}=943\) 尚未全部完成相邻阈值不可行证明。用于生成 \texttt{result4.xlsx} 的同一候选方案统计如下。

\begin{table}[H]
\centering
\caption{问题四当前条件可行方案的类别构成}
\label{tab:q4-candidate}
\begin{tabular}{lrrr}
\toprule
类别 & 保留 & 调整 & 撤销 \\
\midrule
A & 0 & 20 & 0 \\
B & 0 & 39 & 1 \\
C & 5 & 83 & 2 \\
\midrule
合计 & 5 & 142 & 3 \\
\bottomrule
\end{tabular}
\end{table}

该工作簿撤销 B029、C035、C075，共有 147 个活动计划。独立回读重新计算得到 16\,356 个不同占用资源格，连续冲突数、离散冲突数和越界数均为 0，最大资源格占用数为 1，工作簿报告值与重算值一致。表 \ref{tab:q4-candidate} 只描述当前条件可行候选，不把 142 或 943 写成次级全局最优。

\subsubsection{边界敏感性与问题四小结}

在 \(H=638,643,648\) 三个时间域右端点下，均找到通过独立验证的 \(R\le3\) 候选，候选动作数分别为 6\,040、6\,103 和 6\,116，资源格约束数分别为 53\,186、53\,679 和 53\,767。该检验支持“3 台撤销可行上界对边界扰动稳定”，不宣称两个敏感性域的最小撤销数也已经闭合。

问题四在 Q2 的动作框架上只增加 C 类间隔调整，间隔变化通过完整重复事件递推进入资源格约束。六分支不可行证明与三撤销可行见证闭合了当前有限模型下的最小撤销层；其后的类别撤销、调整数量和末级变化量仍需与证据状态一起表述。
